\documentclass[10pt,a4paper]{amsart}
\usepackage[margin=19mm]{geometry}
\usepackage{amsmath,amssymb,mathtools}
\usepackage[scr=rsfs,cal=euler]{mathalfa}
\usepackage{xfrac}
\usepackage[unicode,hidelinks,bookmarks=false]{hyperref}
\newcommand{\ie}{i.e.\ }
\newcommand{\cf}{cf.\ }
\newcommand{\resp}{respectively\ }
\newcommand{\Subsection}{\subsection}
\providecommand{\nobreakdash}{\nobreak\hbox{-}}
\setlength{\emergencystretch}{2em}
\multlinegap0pt
\newcommand{\mr}{\mathscr}
\usepackage{bm}
\usepackage{amssymb}
\usepackage{xcolor}
\usepackage{enumerate}
\newtheorem{prop}{Proposition}
\newtheorem{thm}{Theorem}
\newcommand{\al}{\alpha}
\newcommand{\be}{\beta}
\newcommand{\beqa}{\begin{eqnarray*}}
\newcommand{\cc}{\circ}
\newcommand{\ci}{\subseteq}
\newcommand{\eeqa}{\end{eqnarray*}}
\newcommand{\ga}{\gamma}
\newcommand{\mb}{\mathbb}
\newcommand{\ra}{\rightarrow}
\title{A study on $\mr{F}$-simultaneous approximative $\tau$-compactness property in Banach spaces}
\author{Syamantak Das, Tanmoy Paul}
\date{Source: arXiv:2501.03347v1 (CC BY 4.0)}
\begin{document}
\maketitle

\noindent\emph{Source and licence.} A study on $\mr{F}$-simultaneous approximative $\tau$-compactness property in Banach spaces. Version arXiv:2501.03347v1 (CC BY 4.0), distributed by arXiv under the Creative Commons Attribution 4.0 International licence, http://creativecommons.org/licenses/by/4.0/, as recorded in the arXiv licence metadata for this exact version and shown on the abstract page https://arxiv.org/abs/2501.03347v1. Full licence notice: \texttt{licenses/arxiv-2501.03347-CC-BY-4.0.txt}.\par\smallskip

\noindent\textit{Editorial scope.} Counted unit: the article's thm T9 (source0.tex lines 387--389). The article's complete original proof of it is delivered with it (source0.tex lines 390--417, one \texttt{proof} environment of 28 lines). No mathematics is written, repaired or re-derived: the statement and the proof are the article's own lines, in the article's own notation, and no retained line is substituted or re-flowed.  Retained uncounted as the source-local context the proof needs: lines 375--377.  Nothing was omitted from any retained span, so the package carries no omission record.  No retained line contains a citation, so the wrapper carries no bibliography and no bibliography entry.  No retained line contains a figure environment, an image-inclusion command or drawing code, so no image is delivered and the package declares no asset dependency.  Editorial formatting only, in addition: an AMS \texttt{amsart} preamble that restates the article's own declarations and macros used by the retained lines, namely the environments \texttt{prop}, \texttt{thm} on separate counters, together with the standard packages those lines need.  The retained environments print in this document as Proposition 1, Theorem 1 in order of appearance, whereas the article prints the same items with the numbering of its own full document.  The complete native source of the article as distributed in the arXiv e-print for this exact version is bundled unchanged as \texttt{sources/171231/source0.tex}.
\begin{prop}\label{P5}
For an $N\in\mb{N}$ and monotone norm $\ga:\mb{R}^N\ra\mb{R}_{\geq0}$, if $(f_1,\cdots,f_N)\in NA(\ell_\ga^N(X))$, then $\frac{f_i}{\|f_i\|}\in NA(X)$ for all $i=1,\cdots,N$ such that $f_i\neq0$.
\end{prop}

\begin{thm}\label{T9}
Let $X$ be a $\tau$-ALUR space. Then for every $N\in\mb{N}$, and strictly convex  monotone norm $\ga:\mb{R}^N\ra\mb{R}_{\geq0}$, $\ell_\ga^N(X)$ is a $\tau$-ALUR space.
\end{thm}

\begin{proof}
We consider the case when $N=2$ because the proof for other values of $N$ is similar. Let $\ga:\mb{R}^2\ra\mb{R}_{\geq0}$ be a strictly convex monotone norm and $\ga^\cc$ be the dual norm of $\ga$. Let $(f_1,f_2)\in NA(\ell_\ga^2(X))$ and $((x_n,y_n))\ci S_{\ell_\ga^2(X)}$ be such that $f_1(x_n)+f_2(y_n)\ra 1$.
Let $(z_1,z_2)\in S_{\ell_\ga^2(X)}$ be such that $f_1(z_1)+f_2(z_2)=1$.
 Then
\beqa
1=f_1(z_1)+f_2(z_2)&\leq&\|f_1\|\|z_1\|+\|f_2\|\|z_2\|\\
&\leq&\ga^\cc(\|f_1\|,\|f_2\|)\ga(\|z_1\|,\|z_2\|)=1.
\eeqa
Thus $\|f_1\|\|z_1\|+\|f_2\|\|z_2\|=1$. Thus $(\|f_1\|,\|f_2\|)$ attains its norm at $(\|z_1\|,\|z_2\|)$. 

{\sc Claim: } $(\|x_n\|), (\|y_n\|)$ converges to $\|z_1\|$ and $\|z_2\|$ respectively.  

Let $(\|x_{n_k}\|)$ be a subsequence of $(\|x_n\|)$. Since $(\|x_{n_k}\|)$ is bounded, $(\|x_{n_k}\|)$ has a convergent subsequence, also denoted by $(\|x_{n_k}\|)$, converging to some $\al\in\mb{R}$. In a similar way, we can show that $(\|y_{n_k}\|)$ has a convergent subsequence, also denoted by $(\|y_{n_k}\|)$, converging to some $\be\in\mb{R}$. Since the norm function is continuous, $\ga(\al,\be)=\underset{k}{\lim}\ga(\|x_{n_k}\|,\|y_{n_k}\|)=1$. Now $\|f_1\|\al+\|f_2\|\be=\underset{k}{\lim}(\|f_1\|\|x_{n_k}\|+\|f_2\|\|y_{n_k}\|)=1$. Now by the strict convexity of the norm $\ga$, we have $\al=\|z_1\|$ and $\be=\|z_2\|$. Thus $(\|x_n\|)$ is convergent to $\|z_1\|$ and $(\|y_n\|)$ is convergent to $\|z_2\|$.

{\sc Case 1: } Suppose $f_1,f_2\neq0$. Now 
\beqa 
1=\underset{n}{\lim}(f_1(x_n)+f_2(y_n))&\leq&\underset{n}{\lim\inf}(\|f_1\|\|x_n\|+\|f_2\|\|y_n\|)\\
&\leq&\underset{n}{\lim\sup}(\|f_1\|\|x_n\|+\|f_2\|\|y_n\|)\\
&\leq&\underset{n}{\lim\sup}\ga^\cc(\|f_1\|,\|f_2\|)\ga(\|x_n\|,\|y_n\|)\leq1.
\eeqa
Hence $\underset{n}{\lim}(\|f_1\|\|x_n\|+\|f_2\|\|y_n\|)=\underset{n}{\lim}(f_1(x_n)+f_2(y_n))=1$. Now it is easy to see that $\underset{n}{\lim}(\|f_1\|\|x_n\|-f_1(x_n))=0$ and $\underset{n}{\lim}(\|f_2\|\|y_n\|-f_2(y_n))=0$. Thus $\frac{f_1}{\|f_1\|}(\frac{x_n}{\|x_n\|})\ra 1$ and $\frac{f_2}{\|f_2\|}(\frac{y_n}{\|y_n\|})\ra1$.

From lemma \ref{P5}, $\frac{f_1}{\|f_1\|},\frac{f_2}{\|f_2\|}\in NA(X)$, and so, by our assumption, $(\frac{x_n}{\|x_n\|}),(\frac{y_n}{\|y_n\|})$ are $\tau$-convergent in $X$ and hence $((x_n,y_n))$ is $\tau$-convergent in $\ell_\ga^2(X)$. 



{\sc Case 2:} Suppose $f_1\neq0$ and $f_2=0$. Now $1=\|f_1\|\|x_1\|=\ga^\cc(\|f_1\|,0)\ga(\|x_1\|,0)\leq \ga(\|x_1\|,\|x_2\|)=1$ and hence $\ga(\|x_1\|,0)=1$. Since $(\|f_1\|,0)$ attains its norm on $(\|x_1\|,0)$ and $\ga$ is strictly convex, we have $\|x_2\|=0$ and thus $x_2=0$. Hence $(y_n)$ converges to 0. Also, in a similar way, as in Case 1, $(x_n)$ is $\tau$-convergent. This concludes our assertion. 
\end{proof}

\end{document}
